\section{Introducción}
El presente informe tiene por objetivo el diseño de las fundaciones del edificio en estudio, a partir de las solicitaciones estáticas y sísmicas obtenidas del modelo estructural. El análisis se enmarca en el proyecto de un edificio de hormigón armado de 21 pisos más un nivel de subterráneo, ubicado en Antofagasta, zona sísmica 3, suelo tipo A, con categoría de ocupación II. La estructura se funda sobre 30 elementos verticales: 21 muros de hormigón armado y 9 columnas, que descargan al nivel de fundación las cargas gravitacionales y sísmicas de la superestructura.

En primer lugar, se establecen las bases de diseño materiales, parámetros geotécnicos del sector asignado y criterios geométricos de las zapatas, junto con los esfuerzos de diseño extraídos del modelo y las combinaciones de carga, estáticas y sísmicas, con que se verifica cada fundación. Los muros se resuelven mediante zapatas corridas, verificadas en una dirección, mientras que las columnas se resuelven mediante zapatas aisladas, verificadas de forma biaxial.

En segundo lugar, se desarrolla el diseño de cada fundación, partiendo de la metodología empleada presión de contacto, corte unidireccional, punzonamiento, estabilidad al deslizamiento y armadura de flexión y presentando los resultados del dimensionamiento en planta, la altura adoptada, el porcentaje de compresión y la armadura resultante de los 30 elementos. Las verificaciones se realizan conforme a ACI 318-19, a partir de las solicitaciones sísmicas determinadas según NCh433 y DS 60.